\documentclass[12pt]{article}
\usepackage[utf8]{inputenc}
\usepackage[T1]{fontenc}
\usepackage{hyperref}
\title{Sou Upaa Language Resources for K-12 Curriculum Development}
\author{Compiled by Claude Sonnet 5.5 (claude-sonnet-5-5) for Andr\'as Kornai}
\date{}
\begin{document}
\maketitle

\section{Introduction}
Sou Upaa (wha) is also called Manusela or Wahinama. It is an Austronesian language of the Central Maluku group, spoken in villages in the northern Seram mountains and south toward Teluti Bay, Maluku, Indonesia. Wikipedia gives 7,000 speakers (1989), and Joshua Project gives 7,600 (Ethnologue 2016). The ethnic group, Manusela or Wahai, numbers over 10,100. Note that ``Wahai'' alone can also point to the Saleman language. Very few educational resources were found.

\section{Primary Educational Resources}
No primer, textbook or school programme was found.

\section{Linguistic and Reference Materials}
\begin{itemize}
\item Wikipedia, Manusela language (phonology, classification): \url{https://en.wikipedia.org/wiki/Manusela_language}
\item Numeral system survey page (Sou Upaa): \url{https://lingweb.eva.mpg.de/channumerals/Sou-Upaa.htm}
\end{itemize}
No dictionary or grammar was found.

\section{Religious and Community Resources}
\begin{itemize}
\item Joshua Project profile (Bible translation status: translation started; no resources reported): \url{https://joshuaproject.net/people_groups/15817}
\item Wikipedia, Manusela people: \url{https://en.wikipedia.org/wiki/Manusela_people}
\end{itemize}

\section{Digital and Audio Resources}
No audio or digital resources were found.

\section{Curriculum Development Context}
The language has a sizeable speaker base and is spoken in all generations, but no teaching materials were found and it is not known to be taught in schools. Documentation (a dictionary, an orthography and recordings) would have to come first, done together with the community. Indonesian would likely be the medium for any bridging material.
\end{document}